\documentclass[11pt,a4paper]{article}
\usepackage[T1]{fontenc}
\usepackage{lmodern,amsmath,amssymb,amsthm,microtype}
\usepackage[margin=26mm]{geometry}
\usepackage[colorlinks=true,linkcolor=blue,urlcolor=blue]{hyperref}
\newtheorem{proposition}{Proposition}
\newcommand{\Q}{\mathbb Q}
\newcommand{\Z}{\mathbb Z}
\title{A perturbation obstruction for square-rescaled detector weights}
\author{}\date{17 November 2026}
\begin{document}\maketitle

\begin{abstract}
A balanced detector with a weight whose square class is independent of
the external rational input cannot isolate the integers near a strictly
prepared witness. Its pole preparation is an open condition at finitely
many places. This remains true if the retained elliptic indices satisfy
an independently imposed divisibility condition. The obstruction does
not apply to weights with genuinely input-dependent odd prime support.
\end{abstract}

Fix a nonzero rational $u_0$, a nonzero rational function
$g(X)\in\Q(X)$, and finitely many rational functions
$C_i(X)\in\Q(X)$. These functions may result from fixing every
retained point and all other free coordinates of a proposed source.
Set
\[
 u(X)=u_0g(X)^2,\qquad w=\frac{u}{(1+u)^2},\qquad
 D_i=\frac{C_i^2}{1+uC_i^2},\qquad
 W=\prod_i\left(1+\frac{D_i^2}{u}\right).
\]
Either use one rational slot and set $V=W^{-1}$, or use two rational
slots of this form and set $V=W_1^{-1}+W_2^{-1}$. Put
\[
                         b=\frac{4w}{V^2}.
\]
All displayed divisions are restricted to their nonzero domain.

\begin{proposition}\label{prop:open}
Suppose $X_0\in\Z$ is in this domain, $u(X_0)>0$, and
\[
 0<b(X_0)<1/7,\qquad b(X_0)\in\Q_2^{\times2},\qquad
 v_2(b(X_0))\ge4,\qquad v_7(b(X_0))\ge0.
\]
Suppose also that $b(X_0)$ has square reduced denominator.
There is a real neighborhood $I$ of $X_0$ and a finite collection
of $p$-adic neighborhoods $U_p$ of $X_0$ such that every rational
\[
                       X\in I\cap\bigcap_p U_p
\]
in the nonzero domain has all the same displayed preparation
properties and square reduced denominator. This set contains
infinitely many nonintegers arbitrarily close to $X_0$ in every
specified topology.
\end{proposition}

\begin{proof}
Let $S$ consist of $2,7$ and the primes at which $v_p(u_0)$ is odd.
For $p\notin S$, the valuation of $u(X)$ is even at every rational
input where it is nonzero. The formula
\[
 v_p(b)=v_p(4)+v_p(u)-2v_p(1+u)-2v_p(V)
\]
therefore shows that $v_p(b)$ is even, independently of the tests.
In particular no such prime can give an odd denominator exponent.

At each $p\in S$, rational functions are continuous near $X_0$.
Their nonzero values have locally constant valuations. Shrink a
$p$-adic neighborhood so that all required nonzero denominators
remain nonzero and $v_p(b)$ is unchanged. At two, shrink further
to preserve the nonzero square class; the group
$\Q_2^{\times2}$ is open in $\Q_2^\times$.
The hypotheses at two and seven and the absence of negative odd
valuations at the remaining places of $S$ are preserved.
Real continuity preserves positivity and the strict upper bound.

To see the final assertion directly, choose a positive integer
$N$ having arbitrarily high valuations at the finitely many
specified primes. For primes $q$ outside that finite set and not
dividing $N$, the inputs $X=X_0+N/q$ are nonintegers. As $q$
tends to infinity, they approach $X_0$ in the real topology and
remain in all the prescribed $p$-adic neighborhoods.
\end{proof}

The same argument permits any additional finite collection of
nonzero rational guards and open real or $p$-adic conditions.
It also permits a finite number of independent pole blocks of
the displayed kind. At a rational point of a finite point
algebra, the statement applies to its selected rational
character. Whole-character nonzero guards remain nonzero after
shrinking the neighborhoods, since there are only finitely many
characters.

The intrinsic pole theorem in \cite[Theorem 2.1]{Intrinsic}
supplies rational witnesses for each prepared value $7b(X)$.
Its density statement permits avoidance of any finite collection
of nonzero polynomial norm guards. Thus, when no other closed
equation depending on $X$ is imposed, such pole blocks do not
restore discreteness of the external input.

\paragraph{Elliptic denominator application.}
On a rational square-abscissa chart $x(H)=q^2$, the tentative
weight $u=1/q$ has positive odd valuations precisely at the
odd denominator valuations of $q$. A balanced detector can use
these places to test denominator support inclusion for another
point $U$. Even if that inclusion independently proves
$h\mid j$ for $H=[h]P$, $U=[j]P$, fixing $H,U$ and perturbing
the external input as above leaves these indices unchanged.
Consequently local tests approximating $j/h$ do not alone
prove equality to the external input. Multiplying the weight by
a rational square depending on the input does not change this
conclusion. A new closed source relation or genuinely varying
odd prime support is required.

\begin{thebibliography}{9}
\bibitem{Intrinsic}\emph{A two-coordinate pole test of coefficient
height four}, companion paper,
\nolinkurl{intrinsic_pole_20261117_v56.tex}, Theorem 2.1.
\end{thebibliography}
\end{document}
